\scanpage{209}
\mspage{125}
% archive sheet number: 208
% element: text | top, lines 1-2 | stricte Apparemment
stricte ! Apparemment, les ``majorations''
% element: display | top, line 2 | M pi inclusions
\[
M_\pi\subset\mathrm{Filt}^{\pi^{\pi_0}},\qquad M_{\pi'}\subset\mathrm{Filt}^{\pi'{}^{\pi_0}}
\]
% element: text | upper third, lines 3-5 | sont trop grossières
\unsure{sont trop} grossières. Aussi il s'agit de mieux appréhender les catégories $M_\pi$, pour $\pi\geq\pi_0$\textsuperscript{*}. Je vais hasarder une
% element: margin note | left edge, full height, written sideways | Mais on peut utiliser
% margin note written sideways and slanted
\margin{Mais on peut utiliser 4.5.10, 4.5.11. On en déduit que pour $x\in\mathrm{Ob}\,M_{\pi'}$ \deleted{$\unsure{\operatorname{card}M_\pi|\operatorname{card}M_{\pi'}}$} \inserted{$\subset\mathrm{Ob}\,\mathrm{Filt}^{\pi'{}^{\pi_0}}$}, $\operatorname{card}\mathrm{Fl}(C_0/x)\leq(\pi'{}^{\pi_0})^{\pi_0}=\pi'{}^{\pi_0}$. On \unsure{aura|a} donc $x=\varinjlim_{\unsure{J|I}}x_j$ ($x_j\in\mathrm{Ob}\,C$, $\varinjlim$ filtrante, $\operatorname{card}\mathrm{Fl}J\leq\pi'{}^{\pi_0}$), et \unsure{introduisons|on introduit} $I=\mathrm{ens.}$ des sous-catégories $J'\subset J$ \unsure{telles que} $\operatorname{card}J'\leq\pi$, et puis, \unsure{si} $x_{J'}=\varinjlim_{j\in J'}x_j$, \unsure{on a} $x=\varinjlim_{J'\in I}x_{J'}$, les $x_{J'}\in\mathrm{Ob}\,M_\pi$, $I$ \unsure{grd|grand} devant $\pi$, et $\operatorname{card}I\leq(\operatorname{card}J)^\pi\leq(\pi'{}^{\pi_0})^\pi=\pi'{}^{\pi\pi_0}=\pi'{}^\pi$, \unsure{ça suffit !|ce qui suffit !}}
% element: text | upper middle, proposition heading | Proposition 4.5.13
\uline{Proposition} 4.5.13. Les hypothèses et notations sont celles de 4.5.9.
% element: deletion | upper middle, parenthesis | dont la démonstration
\deleted{(dont la démonstration \unsure{restait|reste|est} incomplète, pour la majoration de cardinaux dans la condition c) des filtrations cardinales).}
% element: text | middle, introduction to conditions | Soit pi un cardinal
Soit $\pi$ un cardinal $\geq\pi_0$, et soit $x\in\mathrm{Ob}\,M$. Considérons les conditions
% element: text | middle, condition a | Il existe une catégorie
\begin{enumerate}[label={},leftmargin=2em]
\item[a)] Il existe une catégorie $I$, avec $\operatorname{card}\mathrm{Fl}I\leq\pi$, et un foncteur $\unsure{\varphi|\phi}:I\longrightarrow C$, \inserted{$i\longmapsto x_i$}, tel que $x$ soit isomorphe à $\varinjlim x_i$.
% element: text | middle, condition a prime | Comme a
\item[a')] Comme a), avec $I$ filtrante.
% element: text | middle, condition a double prime | Comme a
\item[a'')] Comme a), avec $I$ ordonnée filtrante.
% element: insertion | between a double prime and b prime | exists foncteur
\item[b)] \inserted{$\exists$ foncteur \inserted{cofinal} $I\longrightarrow C/x$, avec $\operatorname{card}\mathrm{Fl}I\leq\pi$.}
% element: text | lower middle, condition b prime | La catégorie filtrante
\item[b')] La catégorie filtrante $C/x$ est ``accessible en cardinal $\leq\pi$'', i.e. il existe une catégorie \unsure{ordonnée} \uline{filtrante} $I$, \deleted{avec} \inserted{$\operatorname{card}\mathrm{Fl}I$} $\leq\pi$, et un foncteur cofinal $I\longrightarrow C/x$.
% element: text | lower third, condition b double prime | Comme b prime
\item[b'')] Comme b'), avec de plus $I'$ \unsure{ordonnée}.
% element: text | lower third, condition c | Il existe une famille
\item[c)] Il existe une famille de flèches $(\unsure{u_\alpha|f_\alpha})_{\alpha\in A}$, $x_\alpha\longrightarrow x$, \unsure{épimorphique strict|épimorphique stricte|épimorphiques stricte}, indexée par un ens $A$ tel que $\operatorname{card}A\leq\pi$.
% element: text | foot, condition d | On a x
\item[d)] On a $x\in\mathrm{Ob}\,M_\pi$, i.e. $x$ est $\pi$-accessible dans $M$.
\end{enumerate}
